% !TeX root = ../main.tex
\section{Experiments}

\subsection{Experiments Setting}

\paragraph{Datasets.}
To comprehensively evaluate the effectiveness of our method for multimodal continual test-time adaptation, we conduct experiments on two widely used audio-visual corruption benchmarks, namely Kinetics50-C and VGGSound-C.
Following prior multimodal TTA settings, we apply 15 corruptions to the video modality and 6 corruptions to the audio modality, each with five severity levels.
Detailed construction procedures and protocol descriptions are provided in the appendix.

\paragraph{Baselines.}
We compare ZOSA with representative CTTA methods, including Tent, SAR, READ, ContinualMAE, DeYO, and REM, as well as multimodal adaptation methods including PTA, BECoTTA-MoE, SuMi, and VIDA.
We further include backpropagation-free methods, FOA and T3A.
All baselines are evaluated under the same data streams and protocols, with detailed settings provided in the appendix.


\paragraph{Implementation details.}
For synthetic multimodal domain shifts, we use a CAV-MAE based pretrained audio-visual classifier and follow the standard setup of prior multimodal TTA works.
Unlike conventional batch-based online adaptation, we evaluate all methods in a single-sample streaming setting: test samples arrive sequentially, and the model adapts using only the current sample and a lightweight local history window.
We use ZOO as the online optimizer and update only the selected shift-sensitive layers in the modality branches, while keeping the fusion module frozen.
Unless otherwise specified, we set the streaming window length to $W=24$, select $K=3$ layers for each modality branch, use one ZOO step per test sample, and set the perturbation radius and learning rate to $\delta=10^{-3}$ and $\eta=10^{-2}$, respectively.
The alignment coefficient is set to $\lambda=0.4$.
All experiments are run on a single NVIDIA RTX 4090 GPU and repeated with five random seeds.
We report the average accuracy, with additional implementation details and hyperparameter settings provided in the appendix.


\subsection{Comparision with SOTAs}
We compare against a frozen \texttt{no\_adapt} model and a broad set of continual test-time adaptation baselines, including FOA, T3A, Tent, SAR, READ, ContinualMAE, DeYO, REM, PTA, BECoTTA-MoE, SuMi, and VIDA. Some methods have incomplete verified exports in a subset of co-shift settings. When only protocol-level summaries are verified, we keep the aggregate result in the cross-protocol summary table but omit unverified detailed rows from protocol-specific tables.

\definecolor{zosaBestRow}{HTML}{F8E1CF}
\definecolor{zosaAvgCol}{HTML}{ECEAF8}
\newcommand{\avgcell}[1]{\cellcolor{zosaAvgCol}#1}

\begin{table}[t]
\caption{Kinetics50-C interleaved single-modality corruption. Columns labeled \texttt{V:*} indicate that the audio modality remains clean while only the video modality is corrupted, whereas columns labeled \texttt{A:*} indicate that the video modality remains clean while only the audio modality is corrupted. Accuracy is reported in \%. Higher is better.}
\label{tab:k50_interleaved}
\centering
\scriptsize
\setlength{\tabcolsep}{1pt}
\resizebox{\textwidth}{!}{%
\begin{tabular}{@{}lcccccccccccccccccccccc@{}}
\toprule
Method
& \multicolumn{2}{c}{A clean}
& \multicolumn{1}{c}{V clean}
& \multicolumn{2}{c}{A clean}
& \multicolumn{1}{c}{V clean}
& \multicolumn{2}{c}{A clean}
& \multicolumn{1}{c}{V clean}
& \multicolumn{3}{c}{A clean}
& \multicolumn{1}{c}{V clean}
& \multicolumn{2}{c}{A clean}
& \multicolumn{1}{c}{V clean}
& \multicolumn{2}{c}{A clean}
& \multicolumn{1}{c}{V clean}
& \multicolumn{2}{c}{A clean}
& \cellcolor{zosaAvgCol}Avg. \\
\cmidrule(lr){2-3}\cmidrule(lr){4-4}\cmidrule(lr){5-6}\cmidrule(lr){7-7}
\cmidrule(lr){8-9}\cmidrule(lr){10-10}\cmidrule(lr){11-13}\cmidrule(lr){14-14}
\cmidrule(lr){15-16}\cmidrule(lr){17-17}\cmidrule(lr){18-19}\cmidrule(lr){20-20}
\cmidrule(lr){21-22}
& V:Gaus. & V:Shot & A:Gaus. & V:Imp. & V:Def. & A:Traf. & V:Glas. & V:Mot. & A:Crowd
& V:Zoom & V:Snow & V:Frost & A:Rain & V:Fog & V:Brit. & A:Thun. & V:Contr. & V:Elas. & A:Wind & V:Pix. & V:JPEG & \cellcolor{zosaAvgCol}Avg. \\
\midrule
\multicolumn{23}{c}{$t\ \xrightarrow{\hspace*{18em}}$} \\
\midrule
\texttt{no\_adapt}
& 46.88 & 48.44 & 68.75 & 46.88 & 65.62 & 64.84 & 72.66 & 80.47 & 67.97 & 57.03 & 61.72 & 61.72 & 67.97 & 49.22 & 75.78 & 74.22 & 50.78 & 61.72 & 75.78 & 68.75 & 65.62 & \avgcell{63.47} \\
FOA
& 48.44 & 50.78 & 68.75 & 42.19 & 67.97 & 64.84 & 74.22 & 80.47 & 70.31 & 57.03 & 63.28 & 66.41 & 69.53 & 60.94 & 75.00 & 75.00 & 53.91 & 65.62 & 78.12 & 70.31 & 69.53 & \avgcell{65.36} \\
T3A
& 30.47 & 37.50 & 50.00 & 35.16 & 55.47 & 47.66 & 58.59 & 67.19 & 59.38 & 49.22 & 39.84 & 45.31 & 46.88 & 37.50 & 57.03 & 57.81 & 33.59 & 58.59 & 64.06 & 56.25 & 56.25 & \avgcell{49.70} \\
% ABPEM
% & 39.84 & 41.41 & 57.81 & 42.97 & 57.81 & 60.16 & 54.69 & 62.50 & 59.38 & 50.00 & 46.88 & 54.69 & 59.38 & 53.12 & 60.16 & 67.19 & 37.50 & 59.38 & 64.84 & 53.12 & 57.03 & 54.28 \\
Tent
& 47.66 & 48.44 & 69.53 & 46.88 & 65.62 & 66.41 & 71.88 & 80.47 & 70.31 & 57.03 & 61.72 & 61.72 & 68.75 & 49.22 & 75.00 & 74.22 & 50.78 & 59.38 & 78.12 & 67.19 & 65.62 & \avgcell{63.62} \\
SAR
& 42.97 & 34.38 & 46.88 & 32.03 & 34.38 & 35.16 & 59.38 & 62.50 & 27.34 & 47.66 & 44.53 & 35.16 & 47.66 & 50.00 & 57.81 & 26.56 & 28.91 & 31.25 & 51.56 & 39.06 & 38.28 & \avgcell{41.59} \\
READ
& 46.88 & 49.22 & 68.75 & 46.88 & 65.62 & 66.41 & 70.31 & 79.69 & 70.31 & 58.59 & 60.94 & 65.62 & 68.75 & 53.91 & 75.00 & 73.44 & 51.56 & 61.72 & 76.56 & 68.75 & 67.19 & \avgcell{64.10} \\
ContinualMAE
& 47.66 & 48.44 & 68.75 & 46.88 & 65.62 & 64.84 & 72.66 & 80.47 & 69.53 & 57.03 & 61.72 & 61.72 & 67.97 & 49.22 & 75.78 & 74.22 & 50.78 & 61.72 & 76.56 & 68.75 & 65.62 & \avgcell{63.62} \\
DeYO
& 46.88 & 48.44 & 68.75 & 46.88 & 65.62 & 66.41 & 72.66 & 80.47 & 70.31 & 58.59 & 60.94 & 60.16 & 69.53 & 46.09 & 75.78 & 75.00 & 48.44 & 60.16 & 78.91 & 70.31 & 66.41 & \avgcell{63.65} \\
REM
& 47.66 & 48.44 & 68.75 & 46.88 & 65.62 & 65.62 & 72.66 & 80.47 & 69.53 & 57.03 & 61.72 & 62.50 & 67.97 & 51.56 & 75.00 & 75.00 & 50.78 & 60.94 & 78.12 & 66.41 & 65.62 & \avgcell{63.73} \\
PTA
& 46.88 & 49.22 & 68.75 & 46.09 & 66.41 & 66.41 & 71.88 & 79.69 & 70.31 & 57.03 & 61.72 & 64.84 & 69.53 & 45.31 & 75.00 & 74.22 & 50.00 & 60.94 & 78.91 & 68.75 & 63.28 & \avgcell{63.58} \\
BECoTTA-MoE
& 45.31 & 47.66 & 71.09 & 46.09 & 66.41 & 66.41 & 73.44 & 80.47 & 67.97 & 54.69 & 60.16 & 57.81 & 68.75 & 46.09 & 74.22 & 71.88 & 50.78 & 60.16 & 74.22 & 66.41 & 65.62 & \avgcell{63.54} \\
SuMi
& 47.66 & 49.22 & 69.53 & 46.88 & 67.19 & 65.62 & 72.66 & 79.69 & 68.75 & 53.91 & 62.50 & 65.62 & 67.97 & 55.47 & 74.22 & 71.88 & 50.00 & 63.28 & 76.56 & 69.53 & 64.84 & \avgcell{63.95} \\
VIDA
& 47.66 & 48.44 & 69.53 & 42.97 & 58.59 & 63.28 & 60.16 & 70.31 & 58.59 & 48.44 & 53.91 & 65.62 & 68.75 & 36.72 & 71.09 & 72.66 & 48.44 & 58.59 & 75.78 & 69.53 & 64.06 & \avgcell{59.67} \\
\midrule
\rowcolor{zosaBestRow}
\textbf{ZOSA}
& \textbf{53.91} & \textbf{56.25} & 69.53 & \textbf{46.09} & \textbf{67.19} & 69.53 & \textbf{76.56} & \textbf{81.25} & \textbf{73.44} & \textbf{57.81} & 60.94 & 64.06 & 68.75 & \textbf{66.41} & 71.88 & \textbf{77.34} & 48.44 & 65.62 & \textbf{78.12} & \textbf{75.78} & 69.53 & \textbf{66.59} \\
\bottomrule
\end{tabular}%
}
\end{table}

\begin{table}[t]
\caption{Average accuracy (\%) across six multimodal corruption protocols. Higher is better. The proposed method achieves the best results in five of the six settings and remains competitive on the remaining one.}
\label{tab:main_results}
\centering
\small
\begin{tabular}{lcccccc}
\toprule
Method & K50-I & VGG-I & K50-E & K50-H & VGG-E & VGG-H \\
\midrule
\texttt{no\_adapt} & 63.47 & 45.54 & 47.33 & 44.43 & 13.73 & 12.60 \\
FOA & 65.36 & \textbf{48.21} & 52.60 & 50.70 & 19.87 & 17.83 \\
Tent & 63.62 & 45.72 & 46.40 & 42.03 & 9.57 & 8.77 \\
SAR & 41.59 & 24.85 & 29.27 & 27.53 & 7.20 & 7.17 \\
READ & 64.10 & 44.98 & 49.50 & 46.73 & 5.47 & 5.17 \\
ContinualMAE & 63.62 & 45.54 & 47.27 & 44.33 & 13.70 & 12.47 \\
DeYO & 63.65 & 45.80 & 48.13 & 45.50 & 13.83 & 12.57 \\
REM & 63.73 & 45.91 & 47.17 & 44.07 & 11.73 & 10.67 \\
PTA & 63.58 & 43.75 & 46.37 & 35.13 & 5.40 & 4.87 \\
BECoTTA-MoE & 63.54 & 44.08 & 46.43 & 44.73 & 1.37 & 7.53 \\
T3A & 49.70 & 17.67 & 32.50 & 29.93 & 4.07 & 3.57 \\
SuMi & 63.95 & 6.40 & 48.03 & 44.40 & 5.07 & 1.50 \\
VIDA & 59.67 & 38.73 & 30.75 & 23.75 & 3.00 & 3.50 \\
\textbf{ZOSA (ours)} & \textbf{66.59} & 47.58 & \textbf{56.83} & \textbf{55.63} & \textbf{23.77} & \textbf{22.07} \\
\bottomrule
\end{tabular}
\end{table}

\paragraph{Main results.}
Table~\ref{tab:k50_interleaved} gives the detailed per-corruption view in the grouped layout used throughout our benchmark tables. Each column explicitly indicates which modality remains clean and which modality is corrupted, making the alternating single-modality schedule clear. The \textbf{ZOSA} row shows that our method improves over the frozen model across nearly all local corruption groups and remains competitive with the strongest baselines at fine granularity.

Table~\ref{tab:main_results} then summarizes the protocol-level averages. ZOSA consistently improves over the frozen \texttt{no\_adapt} model across all six evaluation protocols. The gains are especially large under paired cross-modal shifts on Kinetics50-C, where our method improves the average accuracy from $47.33\%$ to $56.83\%$ on PairTrack Easy and from $44.43\%$ to $55.63\%$ on PairTrack Async Hard. These margins indicate that explicit multimodal adaptation is particularly valuable when the two modalities drift jointly rather than independently.

On the two interleaved settings, ZOSA remains highly competitive. It reaches $66.59\%$ on Kinetics50-C and $47.58\%$ on VGGSound-C, clearly improving over the frozen baseline and outperforming FOA on Kinetics50-C while remaining within only $0.63$ points on VGGSound-C. The same pattern is even more pronounced on the VGGSound-C co-shift protocols, where ZOSA improves over \texttt{no\_adapt} by roughly ten absolute points and also surpasses the strongest competing FOA variant by $3.90$ and $4.24$ points on PairTrack Easy and Async Hard, respectively.

\paragraph{Discussion.}
The results suggest a clear robustness profile. First, our method is not merely recovering clean-model performance after mild corruption; it remains effective under harder co-shift scenarios where both modalities are perturbed simultaneously. Second, the relative advantage is larger on Kinetics50-C than on VGGSound-C, which suggests that the current design handles structured co-shifts better than highly diverse open-domain audio-visual noise. Third, several classic CTTA baselines such as Tent, SAR, and PTA degrade noticeably on the harder multimodal settings, whereas ZOSA remains stable. This behavior is consistent with the goal of learning a more reliable multimodal adaptation rule instead of relying on single-modality entropy minimization alone.

\paragraph{Reporting note.}
The full per-corruption tables are provided in Appendix~\ref{tab:appendix_k50_interleaved}--\ref{tab:appendix_vgg_hard}. In the main paper, we keep only the representative grouped table and the protocol-level summary because they communicate the core message more clearly and avoid overwhelming the reader with overly wide tables.



\subsection{Ablation Study}

\subsection{Further Analysis}
